\definecolor{decFill}{HTML}{F1EFE8}
\definecolor{decLine}{HTML}{5F5E5A}
\definecolor{linFill}{HTML}{E6F1FB}
\definecolor{linLine}{HTML}{185FA5}
\definecolor{nlFill}{HTML}{FAECE7}
\definecolor{nlLine}{HTML}{993C1D}
\definecolor{anFill}{HTML}{E1F5EE}
\definecolor{anLine}{HTML}{0F6E56}
\definecolor{noteFill}{HTML}{FAEEDA}
\definecolor{noteLine}{HTML}{854F0B}
\definecolor{boxFill}{HTML}{FBFAF7}
\definecolor{boxLine}{HTML}{888780}

\begin{figure}[H]
\centering
\resizebox{\linewidth}{!}{%
\begin{tikzpicture}[
    font=\small,
    >={Stealth[length=2.4mm]},
    node distance=6mm and 12mm, % Reduced from 8mm and 14mm to save vertical space
    basebox/.style={
        rectangle,
        rounded corners=3pt,
        line width=0.4pt,
        align=center,
        inner sep=5pt
    },
    dec/.style={basebox, draw=decLine, fill=decFill, text width=44mm},
    lin/.style={basebox, draw=linLine, fill=linFill, text width=46mm},
    nl/.style ={basebox, draw=nlLine,  fill=nlFill,  text width=46mm},
    an/.style ={basebox, draw=anLine,  fill=anFill,  text width=58mm},
    fbox/.style={
        basebox,
        draw=boxLine,
        fill=boxFill,
        text width=42mm,
        minimum height=9mm
    },
    pbox/.style={
        basebox,
        draw=boxLine,
        fill=boxFill,
        text width=52mm,
        minimum height=9mm
    },
    elabel/.style={font=\scriptsize, inner sep=1.5pt},
    flow/.style={->, draw=decLine, line width=0.5pt}
]

%--------------------------------------------------
% Helper commands
%--------------------------------------------------

\newcommand{\question}[4][]{%
    \ifx\relax#4\relax
        \node[dec,#1] (#2) {#3};%
    \else
        \node[dec,#1] (#2) {#3\\[1pt]\scriptsize #4};%
    \fi
}

\newcommand{\resultbox}[4][]{%
    \node[#2,#1] (#3) {#4};
}

\newcommand{\yesarrow}[2]{%
    \draw[flow] (#1) -- (#2) node[elabel, midway, above] {Yes};
}

\newcommand{\noarrow}[2]{%
    \draw[flow] (#1) -- (#2) node[elabel, midway, left] {No};
}

%--------------------------------------------------
% Main decision tree
%--------------------------------------------------

\question{start}{Start: what to predict?}{}

\question[below=of start]
    {q1}
    {Global response, load path or modal behaviour?}
    {}

\question[below=of q1]
    {q2}
    {Clamped-member stress or fatigue life?}
    {}

\question[below=of q2]
    {q3}
    {Local bolt behaviour?}
    {Fatigue-critical detail, e.g.\ under-head fillet stress}

% Reduced from 42mm to 24mm to heavily compress vertical whitespace
\question[below=24mm of q3]
    {q4}
    {See a phenomenon occur?}
    {To tell the scale of a phenomenon.}

\resultbox[right=of q1]{lin}{r1}
    {Beam or bonded contact}

\resultbox[right=of q2]{lin}{r2}
    {Beam or CBUSH\\[1pt]\scriptsize Beam preferred, see Section \ref{chp:results}}

\question[right=of q3]
    {q3b}
    {Thread-root phenomena?}
    {}

\resultbox[right=of q4]{nl}{r4}
    {Solid bolt with threads\\[1pt]\scriptsize + nonlinear transient analysis}

\resultbox[below left=8mm and -18mm of q3b] % Reduced vertical distance
    {lin}
    {s1}
    {Solid bolt, no threads}

\resultbox[below right=8mm and -22mm of q3b] % Reduced vertical distance
    {lin}
    {s2}
    {Solid bolt with threads}

\node[an, below=6mm of q4] (anode) % Reduced from 10mm
{
Simplified model + analytical checks\\[1pt]
\scriptsize Binary answers
};

%--------------------------------------------------
% Flow arrows
%--------------------------------------------------

\draw[flow] (start) -- (q1);

\noarrow{q1}{q2}
\noarrow{q2}{q3}
\noarrow{q3}{q4}
\noarrow{q4}{anode}

\yesarrow{q1}{r1}
\yesarrow{q2}{r2}
\yesarrow{q3}{q3b}
\yesarrow{q4}{r4}

\draw[flow]
(q3b.south) -- ++(0,-4mm) -| (s1.north)
node[elabel,pos=0.25,above] {No};

\draw[flow]
(q3b.south) -- ++(0,-4mm) -| (s2.north)
node[elabel,pos=0.25,above] {Yes};

%--------------------------------------------------
% Analytical checks
%--------------------------------------------------

\node[fbox,
      below=10mm of anode.south west, % Reduced from 18mm
      anchor=north west,
      xshift=2mm] (f1)
{$F_{shear}<F_M\,\mu$};

\node[pbox,right=8mm of f1] (p1)
{
\textbf{Sliding / self-loosening}\\[1pt]
\scriptsize Slip is the gatekeeper phenomenon
};

\node[fbox,below=2mm of f1] (f2) % Reduced from 4mm
{$F_Z=\dfrac{f_Z}{\delta_S+\delta_P}$};

\node[pbox,right=8mm of f2] (p2)
{
\textbf{Embedding}\\[1pt]
\scriptsize Preload loss from surface smoothing
};

\node[fbox,below=2mm of f2] (f3) % Reduced from 4mm
{$\sigma_{bend}=\dfrac{M}{W}$};

\node[pbox,right=8mm of f3] (p3)
{\textbf{Bolt bending}};

\node[fbox,below=2mm of f3] (f4) % Reduced from 4mm
{$\sigma_{bearing}=\dfrac{F_{shear}}{d_{bolt}\cdot t}$};

\node[pbox,right=8mm of f4] (p4)
{
\textbf{Hole elongation}\\[1pt]
\scriptsize Bearing against the hole wall
};

\node[fbox,below=2mm of f4] (f5) % Reduced from 4mm
{$\tau_M=\dfrac{M_G}{W_P}$};

\node[pbox,right=8mm of f5] (p5)
{\textbf{Residual torsion}};

\node[fbox,below=2mm of f5] (f6) % Reduced from 4mm
{$F_M>(1-\Phi)\,F_L$};

\node[pbox,right=8mm of f6] (p6)
{
\textbf{Joint separation}\\[1pt]
\scriptsize Residual clamp force stays positive
};

%--------------------------------------------------
% Analytical-checks container
%--------------------------------------------------

% FIX: Removed draw and fill properties from this node so it acts 
% as an invisible bounding box, removing the "double line" bug.
\node[
    fit=(f1)(p1)(f6)(p6),
    inner xsep=7pt,
    inner ysep=7pt
] (anfit) {};

\coordinate (hl) at ([yshift=6mm]anfit.north west);
\coordinate (hr) at ([yshift=6mm]anfit.north east);

\begin{scope}[on background layer]
    \node[
        rectangle,
        rounded corners=5pt,
        draw=anLine,
        fill=anFill!35,
        line width=0.4pt,
        fit=(anfit)(hl)(hr),
        inner sep=0pt
    ] (anbox) {};
\end{scope}

\node[
    anchor=north west,
    font=\scriptsize\bfseries,
    text=anLine,
    inner sep=3pt
]
at (anbox.north west)
{Analytical checks};


\end{tikzpicture}
}

\caption{
Decision flowchart for selecting a bolt FE modelling strategy.
Each question is answered \emph{Yes}
(the simpler model on the right is sufficient)
or \emph{No}
(more fidelity is required, continue down).
Beam and CBUSH are capability-equivalent for global and
clamped-member quantities; the beam was the more efficient
choice for the joint studied here, offering greater accuracy
for less preprocessing than a VDI~2230 CBUSH
(Section~\ref{chp:discussion}).
The analytical checks are the subset of phenomena with a
closed-form criterion; the full set is given in the
phenomenon--method matrix
(Table~\ref{tab:phenomenaMethodMatrix}).
}
\label{fig:decisionTree}

\end{figure}